\documentclass[11pt]{article}

\usepackage[margin=1.1in]{geometry}
\usepackage{amsmath,amssymb}
\usepackage{booktabs}
\usepackage{microtype}
\usepackage[hidelinks]{hyperref}

\newcommand{\Rb}{\mathbb{R}}
\newcommand{\dotp}[2]{\left\langle #1,\, #2 \right\rangle}
\newcommand{\ucon}{u_{\mathrm{contrast}}}
\newcommand{\bpre}{b_{\mathrm{pre}}}
\newcommand{\xnorm}{x_{\mathrm{norm}}}

\title{Causal contribution, from the ground up\\[2pt]
       \large Where $\varphi$, $z$ and $\ell_2$ come from}
\author{MrVLA}
\date{}

\begin{document}
\maketitle

\begin{abstract}
\noindent
This note builds the causal-contribution measure from the robot outward, in eight steps.
Nothing is assumed except what a dot product is. The endpoint is the formula
\[
\varphi_j \;=\; \frac{\ell_2}{r}\, z_j \dotp{w_j}{g \odot \ucon},
\]
but the point of the note is that this formula is never \emph{chosen} --- it falls out of
substituting one equation into another. Companion to
\texttt{notes/causal\_contribution.tex}, which states the same result compactly for the paper.
\end{abstract}

\paragraph{Notation.} $\dotp{a}{b}$ is a dot product: multiply matching entries and add them
all up, giving one number. $a \odot b$ is an \emph{entrywise} product: multiply entry~1 by
entry~1, entry~2 by entry~2, and so on, giving a vector back with no summing. The two are
related by a fact used repeatedly below,
\begin{equation}
\dotp{a \odot g}{b} \;=\; \dotp{a}{g \odot b},
\label{eq:move}
\end{equation}
which just says a fixed entrywise scaling can be applied to either side of a dot product.

\section{What the robot is actually doing}
\label{sec:robot}

At every timestep OpenVLA outputs seven numbers: how far to move in $x$, $y$ and $z$, how to
rotate (roll, pitch, yaw), and whether to open or close the gripper.

It does not emit these as raw continuous values. Each one is chosen from \textbf{256 discrete
options}, called \emph{bins}. Bin~0 might mean ``move hard left'', bin~128 ``do not move'',
bin~255 ``move hard right''.

To choose, the model assigns all 256 bins a score and takes the largest:
\begin{center}
\begin{tabular}{@{}rcl@{}}
\toprule
bin & & score \\
\midrule
$0$   & $\rightarrow$ & $1.2$ \\
$1$   & $\rightarrow$ & $0.8$ \\
$\vdots$ & & $\vdots$ \\
$173$ & $\rightarrow$ & $\mathbf{9.4}$ \quad $\leftarrow$ largest, so this is the action \\
$\vdots$ & & $\vdots$ \\
$255$ & $\rightarrow$ & $2.1$ \\
\bottomrule
\end{tabular}
\end{center}

So the question \emph{``what caused this action?''} is really the question \textbf{``what
pushed bin 173's score above all the others?''} Everything that follows is an answer to that
question and nothing more.

\section{Where the scores come from}
\label{sec:scores}

At the final layer the model holds a single vector $h \in \Rb^{d}$ --- roughly $4096$ numbers.
It is the model's complete internal state at the instant of the decision.

Every bin $t$ owns a fixed vector $u_t \in \Rb^{d}$, learned during pre-training and unchanged
at inference. The score for bin $t$ is a dot product of the two, with a normalisation applied
to $h$ first:
\begin{equation}
\operatorname{score}(t) \;=\; \dotp{\tfrac{h}{r} \odot g}{u_t}.
\label{eq:score}
\end{equation}

\begin{center}
\begin{tabular}{@{}lll@{}}
\toprule
symbol & size & changes per decision? \\
\midrule
$h$   & $d$ numbers & yes --- the state right now \\
$r$   & $1$ number  & yes --- a scale factor computed from $h$ \\
$g$   & $d$ numbers & no --- a fixed learned gain, applied entrywise \\
$u_t$ & $d$ numbers & no --- bin $t$'s fixed vector \\
\bottomrule
\end{tabular}
\end{center}

\paragraph{The single most important fact in this note.} Between $h$ and the score there is
\emph{nothing but multiplying and adding}. No ReLU, no attention, no nonlinearity of any kind.

The consequence is what the entire method rests on: \textbf{if you can split $h$ into pieces,
the score splits into matching pieces.} Cut $h$ into three parts and each part contributes its
own share of the score, and those three shares add back to the total \emph{exactly} --- not
approximately. This is why we work at the last layer, and only at the last layer. Anywhere
earlier there are nonlinearities in the way and the split would have to be estimated instead.

\section{The SAE splits \texorpdfstring{$h$}{h} into pieces}
\label{sec:sae}

The sparse autoencoder is a dictionary of $F = 2048$ directions $w_1, \dots, w_{F}$. Each
$w_j$ lives in the same $\Rb^{d}$ as $h$, and each has length exactly $1$. Think of them as
ingredients.

On any single decision only $k = 100$ of the $2048$ are switched on. The SAE's claim is:
\[
h \;\;\approx\;\; (\text{some amount of } w_5) \;+\; (\text{some amount of } w_{88})
   \;+\; \cdots \text{100 terms} \cdots \;+\; \text{leftovers}.
\]

The ``some amount of'' numbers are the code $z$:
\begin{center}
\begin{tabular}{@{}ll@{}}
\toprule
\multicolumn{2}{@{}l}{$z \in \Rb_{\ge 0}^{F}$, one entry per ingredient} \\
\midrule
$z_j$ & how much of ingredient $j$ is present in $h$ right now \\
      & at most $100$ of the $2048$ entries are nonzero \\
      & never negative \\
\bottomrule
\end{tabular}
\end{center}

That is all $z$ is: \textbf{presence}. How much of each ingredient is in the mix.

\section{The size problem, and where \texorpdfstring{$\ell_2$}{l2} comes from}
\label{sec:l2}

Here is the wrinkle. The SAE was never trained on $h$ directly. It was trained only on vectors
of \textbf{length exactly $1$}.

The reason is practical. $h$ is loud on some decisions and quiet on others, and that overall
loudness is not interesting --- we do not want the dictionary spending its capacity describing
volume. So the loudness is stripped off before encoding and added back afterwards. Three steps
(this is \texttt{\_prenorm} in \texttt{mrvla/train\_sae.py}):

\begin{enumerate}
\item Subtract $\bpre$, a fixed learned offset, the same on every decision.
\item Subtract $\mu$, the average of the vector's own entries, so that they now sum to zero.
      Note $\mu$ is \emph{one number}, subtracted from all $d$ entries --- which is why it
      appears as $\mu\mathbf{1}$ later, with $\mathbf{1} \in \Rb^{d}$ the all-ones vector.
\item Divide by the vector's own length. It now has length exactly $1$. Call the result
      $\xnorm$.
\end{enumerate}

\begin{equation}
\boxed{\;\ell_2 \;=\; \bigl\| \, h - \bpre - \mu\mathbf{1} \, \bigr\|_2\;}
\qquad\text{--- the length we divided by in step 3.}
\label{eq:l2}
\end{equation}

One positive number, different on every decision. Loud decision, large $\ell_2$.

To recover the real $h$, undo the three steps: multiply by $\ell_2$ and add the pieces back.
\begin{equation}
h \;=\; \underbrace{\ell_2 \sum_{j=1}^{F} z_j w_j}_{\text{the ingredients, rescaled}}
\;+\; \mu\mathbf{1} \;+\; \bpre \;+\; \underbrace{e}_{\text{error}} .
\label{eq:sae}
\end{equation}

In words: \emph{$h$ is a weighted mix of 100 ingredients, scaled back up to full size, plus a
fixed offset, plus whatever the SAE could not capture.} Equation~\eqref{eq:sae} is the one
everything else is built on.

\subsection*{Two scale factors, easily confused}

\begin{center}
\begin{tabular}{@{}lll@{}}
\toprule
 & comes from & is \\
\midrule
$\ell_2$ & the SAE   & $\|h - \bpre - \mu\mathbf{1}\|_2$, the per-sample normaliser \\
$r$      & the model & $\sqrt{\operatorname{mean}(h^{2}) + \varepsilon}$, the RMSNorm denominator \\
\bottomrule
\end{tabular}
\end{center}

\noindent
Both are one positive number per decision. They come from different places and mean different
things, and both appear in the final formula. Interchanging them is a silent bug.

\section{How \texorpdfstring{$z$}{z} is actually computed}
\label{sec:z}

$z$ is not fitted or solved for per decision. It is one matrix multiply followed by a sort:

\begin{enumerate}
\item Start from $\xnorm$ --- that is $h$ after the three stripping steps, so it has length $1$.
\item Score how well each of the $2048$ ingredients matches it:
      \[
      \mathrm{pre}_j \;=\; \dotp{\mathrm{enc}_j}{\xnorm}, \qquad j = 1, \dots, F,
      \]
      where $\mathrm{enc}_j$ is column $j$ of the learned encoder matrix $W_{\mathrm{enc}}$.
\item Keep the $k = 100$ largest of those scores. Set the other $1948$ to \emph{exactly} zero.
\item Clip anything negative to zero (ReLU). What remains is $z$.
\end{enumerate}

So $z_j$ is nonzero only if ingredient $j$ was among the $100$ best matches on that decision.
The sparsity is \textbf{architectural} --- a hard ``keep 100, delete the rest'' --- rather than
a penalty that merely encourages small values. That is the TopK design of Gao et al.\ (2024),
and it means $k$ is exact and there is no shrinkage bias on the surviving values.

\paragraph{Nothing is trained at attribution time.} The dictionary itself ($W_{\mathrm{enc}}$,
$W_{\mathrm{dec}}$, $\bpre$) is learned once, ahead of time, from collected activations. To
attribute, we run this forward pass on each stored $h$ and record $z$, $\mu$ and $\ell_2$
alongside it.

\section{Deriving the formula}
\label{sec:derivation}

We now have both halves:
\begin{align}
\text{the score:} \qquad & \operatorname{score}(t) = \dotp{\tfrac{h}{r} \odot g}{u_t}
  \tag{\ref{eq:score}} \\
\text{the split:} \qquad & h = \ell_2 \sum_j z_j w_j + \mu\mathbf{1} + \bpre + e
  \tag{\ref{eq:sae}}
\end{align}
Substitute the second into the first, in four steps.

\subsection*{Step A --- only differences between bins matter}

The winner is whichever bin scores highest. Adding the same constant to all $256$ scores does
not change the winner. So the absolute score is not the interesting quantity; the \emph{gap
above the average bin} is.

Define the contrast direction by subtracting the average of all $256$ bin vectors,
\begin{equation}
\ucon \;=\; u_t \;-\; \frac{1}{256}\sum_{s=1}^{256} u_s ,
\label{eq:ucon}
\end{equation}
and define $\Lambda$ as the emitted bin's gap above the average score. Using \eqref{eq:move} to
move $g$ across the dot product,
\begin{equation}
\Lambda \;=\; \operatorname{score}(t) - \frac{1}{256}\sum_{s}\operatorname{score}(s)
        \;=\; \frac{1}{r}\dotp{h}{g \odot \ucon}.
\label{eq:lambda}
\end{equation}

$\Lambda$ is \textbf{the quantity that decides the action}, and it is what we attribute.
Anything that lifts all $256$ bins equally now scores zero, because it cancels in the
subtraction. Only things that pick a \emph{favourite} count. This is why the measure separates
``choice'' from ``confidence'' --- and why a naive attribution to $u_t$ rather than $\ucon$
gives every high-norm direction a score it has not earned.

\subsection*{Step B --- substitute}

Replace $h$ in \eqref{eq:lambda} using \eqref{eq:sae}:
\begin{equation}
\Lambda \;=\; \frac{1}{r}\dotp{\;\ell_2 \sum_j z_j w_j \;+\; \mu\mathbf{1} \;+\; \bpre \;+\; e\;}{\,g \odot \ucon\,}.
\label{eq:subbed}
\end{equation}

\subsection*{Step C --- distribute}

A dot product distributes over addition, $\dotp{a+b}{c} = \dotp{a}{c} + \dotp{b}{c}$. That is
the only real move in the whole derivation. Pulling \eqref{eq:subbed} apart term by term:
\begin{equation}
\Lambda \;=\;
\underbrace{\sum_{j} \frac{\ell_2}{r}\, z_j \dotp{w_j}{g \odot \ucon}}_{\text{the ingredients}}
\;+\;
\underbrace{\frac{1}{r}\dotp{\mu\mathbf{1} + \bpre}{g \odot \ucon}}_{\text{the fixed offset}}
\;+\;
\underbrace{\frac{1}{r}\dotp{e}{g \odot \ucon}}_{\text{the leftovers}} .
\label{eq:decomp}
\end{equation}

\subsection*{Step D --- name it}

\begin{equation}
\boxed{\;\varphi_j \;=\; \frac{\ell_2}{r}\, z_j \dotp{w_j}{g \odot \ucon}\;}
\label{eq:phi}
\end{equation}

That is it. $\varphi_j$ is one term of the first sum in \eqref{eq:decomp}. Nothing was invented
or chosen: it fell out of substituting one equation into another and distributing.

\section{Reading the formula}
\label{sec:reading}

\begin{equation*}
\varphi_j \;=\;
\underbrace{\frac{\ell_2}{r}}_{\text{(1) unit conversion}}
\;\cdot\;
\underbrace{\vphantom{\frac{\ell_2}{r}}\; z_j \;}_{\text{(2) presence}}
\;\cdot\;
\underbrace{\vphantom{\frac{\ell_2}{r}}\; \dotp{w_j}{g \odot \ucon} \;}_{\text{(3) geometry}}
\end{equation*}

\paragraph{(1) $\ell_2/r$ --- unit conversion.} $\ell_2$ converts out of the SAE's
length-one world back to real size; $1/r$ converts from residual units into score units.
Together they only ensure the number means ``score points''.

\paragraph{(2) $z_j$ --- presence.} How much of ingredient $j$ is in the mix right now.
\textbf{This factor, alone, is what a firing-based metric looks at.}

\paragraph{(3) $\dotp{w_j}{g \odot \ucon}$ --- geometry.} Does ingredient $j$'s direction point
towards the bin that won, as against the average bin? This is what makes the measure
\emph{causal}, and it is the part a firing-based metric has no access to.

\medskip
\noindent
The consequence is the argument of the paper in two lines:
\begin{center}
\begin{tabular}{@{}lll@{}}
\toprule
condition & result & name \\
\midrule
$z_j$ large, factor (3) $\approx 0$ & $\varphi_j \approx 0$ &
  \textbf{busy but inert} --- fires constantly, changes nothing \\
$z_j$ rarely on, factor (3) large & $\varphi_j$ large when it fires &
  \textbf{rare but decisive} --- seldom fires, but flips the action \\
\bottomrule
\end{tabular}
\end{center}
\noindent
A firing metric sees only column (2). It therefore calls the first case important and misses
the second entirely.

\section{Why this is checkable rather than merely plausible}
\label{sec:check}

Because Step~C was pure algebra, the three groups in \eqref{eq:decomp} \emph{must} sum to
$\Lambda$. So one can measure what fraction of the decision each group accounts for, and the
three fractions are forced to sum to $1$. Measured on the goal suite at layer 31 with $k=100$:

\begin{center}
\begin{tabular}{@{}lr@{}}
\toprule
group & share of the decision \\
\midrule
ingredients (the features)          & $0.531$ \\
fixed offset ($\mu\mathbf{1} + \bpre$) & $0.405$ \\
leftovers (SAE error $e$)           & $0.064$ \\
\midrule
total                               & $1.000$ \quad \emph{by construction} \\
\bottomrule
\end{tabular}
\end{center}

\noindent
This is the \textbf{sufficiency gate}: features plus offset account for $0.936$ of the
decision, against a threshold of $0.80$ fixed in advance. The middle row is itself a small
finding --- roughly $40\%$ of the action margin is a constant lean toward a rest pose, present
on every decision and belonging to no feature at all.

A separate check confirms the arithmetic itself. Recomputing $\Lambda$ by summing the
$\varphi_j$ and comparing against the true value gives a mean absolute discrepancy of
$2.7 \times 10^{-14}$ --- floating-point dust. The decomposition is not close; it is exact.

\end{document}
